% Generated draft: requires mathematical and source review.
\documentclass[11pt]{article}
\usepackage[margin=25mm]{geometry}
\usepackage{fontspec}
\setmainfont{Latin Modern Roman}
\usepackage{amsmath,amssymb,mathtools,unicode-math}
\setmathfont{Latin Modern Math}
\usepackage{enumitem,needspace,xurl,hyperref}
\setlength{\parindent}{0pt}
\setlength{\parskip}{5pt}
\setlength{\emergencystretch}{3em}
\newcommand{\N}{\mathbb{N}}
\newcommand{\Z}{\mathbb{Z}}
\newcommand{\Q}{\mathbb{Q}}
\newcommand{\R}{\mathbb{R}}
\newcommand{\C}{\mathbb{C}}
\newcommand{\F}{\mathbb{F}}
\newcommand{\A}{\mathbb{A}}
\newcommand{\PP}{\mathbb{P}}
\newcommand{\E}{\mathbb{E}}
\newcommand{\eps}{\varepsilon}
\newcommand{\dd}{\,\mathrm{d}}
\begin{document}
\section*{UnsolvedMath ID: 700578}
\textbf{Source problem code:} catalogue entry
\subsection{Definitions and mathematical statement}
Let \(n\geq 1\) be an integer and let \(a=(a_1,\ldots,a_n)\in\Q^n\subseteq\R^n\). Define
\[
S(a)=\{ka+b:k\in\Z,\ b\in\Z^n\}\subseteq\R^n,
\qquad
\langle x,y\rangle=\sum_{j=1}^n x_jy_j.
\]
As stated in the source, \(S(a)\) is a discrete additive subgroup of \(\R^n\) of rank \(n\). An orthogonal basis of this subgroup means a \(\Z\)-basis \(v_1,\ldots,v_n\in S(a)\) such that
\[
S(a)=\bigoplus_{i=1}^n\Z v_i,
\qquad
v_i\ne 0\quad(1\leq i\leq n),
\qquad
\langle v_i,v_j\rangle=0\quad(i\ne j).
\]
Find necessary and sufficient conditions, expressed in terms of the rational coordinates \(a_1,\ldots,a_n\), for such an orthogonal basis to exist.

The supplied registry labels this problem open. Its accompanying discussion reports an explicitly unverified AI claim that conditions have been found; no such conditions are supplied here, and that claim is not treated as an established solution.
\subsection{Short English statement}
Determine exactly which rational vectors \(a\) make the lattice \(\Z^n+\Z a\) admit a basis over \(\Z\) whose vectors are pairwise orthogonal for the standard Euclidean scalar product.
\subsection{Sources}
\begin{itemize}
\item Supplied local UnsolvedMath registry record: permanent ID \(700578\), exact source code \texttt{catalogue entry}, title \emph{Kourovka Notebook Problem 9.45}.
\item Yu. D. Popov and I. V. Protasov, problem 9.45, \emph{The Kourovka Notebook (Unsolved Problems in Group Theory)}, 9th issue (1984). The supplied record identifies the later source as the 21st issue edition (October 2026), \texttt{arXiv:1401.0300v48}; its collection description names E. I. Khukhro and V. D. Mazurov as editors.
\item Supplied discussion: \emph{[Unverified AI claim]} that conditions were found, attributed to M. Kinyon and J. Urban, \emph{Preprint}, 2026: \url{https://kourovkanotebookorg.wordpress.com/wp-content/uploads/2026/09/kourovka-mathematics.pdf}. The supplied status review says this claim is marked with \(\#\) in the 21st issue and classifies its evidence as \texttt{unrefereed\_claim}.
\item The supplied status assessment is dated \texttt{2026-10-16} and retains the canonical label \texttt{open}. This is supplied dated curation, not an independently verified literature assessment; no cited source was visited for this draft.
\end{itemize}
\end{document}
